\documentclass[10pt,a4paper]{amsart}
\usepackage[margin=19mm]{geometry}
\usepackage{amsmath,amssymb,mathtools}
\usepackage[scr=rsfs,cal=euler]{mathalfa}
\usepackage{xfrac}
\usepackage[unicode,hidelinks,bookmarks=false]{hyperref}
\newcommand{\ie}{i.e.\ }
\newcommand{\cf}{cf.\ }
\newcommand{\resp}{respectively\ }
\newcommand{\Subsection}{\subsection}
\providecommand{\nobreakdash}{\nobreak\hbox{-}}
\setlength{\emergencystretch}{2em}
\multlinegap0pt
\usepackage[shortlabels]{enumitem}
\usepackage{amssymb}
\usepackage{mathtools}
\usepackage{bm}
\usepackage{dsfont}
\usepackage[normalem]{ulem}
\usepackage{xcolor}
\usepackage{tikz}
\newtheorem{hypothesis}{Hypothesis}
\newtheorem{lemma}{Lemma}
\DeclareMathOperator{\supp}{supp}
\title{Selection principles for quasi-stationary distributions and reinforcement processes}
\author{Michel~Bena\"im$^{\dagger}$, Tresnia~Berah$^{\ddagger}$, Pierre~Germain$^{\dagger,\S}$}
\date{Source: arXiv:2606.25857v3 (CC BY 4.0)}
\begin{document}
\maketitle

\noindent\emph{Source and licence.} Selection principles for quasi-stationary distributions and reinforcement processes. Version arXiv:2606.25857v3 (CC BY 4.0), distributed by arXiv under the Creative Commons Attribution 4.0 International licence, http://creativecommons.org/licenses/by/4.0/, as recorded in the arXiv licence metadata for this exact version and shown on the abstract page https://arxiv.org/abs/2606.25857v3. Full licence notice: \texttt{licenses/arxiv-2606.25857-CC-BY-4.0.txt}.\par\smallskip

\noindent\textit{Editorial scope.} Counted unit: the article's lemma lem:block-time-limsup (source0.tex lines 2858--2868). The article's complete original proof of it is delivered with it (source0.tex lines 2870--3232, one \texttt{proof} environment of 363 lines). No mathematics is written, repaired or re-derived: the statement and the proof are the article's own lines, in the article's own notation, and no retained line is substituted or re-flowed.  Retained uncounted as the source-local context the proof needs: lines 572--574; lines 829--868; lines 986--994; lines 2827--2835.  Nothing was omitted from any retained span, so the package carries no omission record.  No retained line contains a citation, so the wrapper carries no bibliography and no bibliography entry.  No retained line contains a figure environment, an image-inclusion command or drawing code, so no image is delivered and the package declares no asset dependency.  Editorial formatting only, in addition: an AMS \texttt{amsart} preamble that restates the article's own declarations and macros used by the retained lines, namely the environments \texttt{hypothesis}, \texttt{lemma} on separate counters, together with the standard packages those lines need.  The retained environments print in this document as Hypothesis 1, Lemma 1 in order of appearance, whereas the article prints the same items with the numbering of its own full document.  The complete native source of the article as distributed in the arXiv e-print for this exact version is bundled unchanged as \texttt{sources/203650/source0.tex}.
\begin{hypothesis}\label{hyp:main}
The cemetery state \(\partial\) is accessible from every \(i\in S\).
\end{hypothesis}

\begin{lemma}[Support properties of a top right eigenvector]
\label{lem:right-perron-support}
There exists a vector \(h\in\mathbb R^S\), with \(h\geq0\) and \(h\neq0\),
such that
\[
Ph=\rho_\kappa h.
\]
For any such vector, let \(U:=\supp(h)\). Then:
\begin{enumerate}[label=\textup{(\roman*)}]
\item \(U\) is a union of communicating classes and \(U^c\) is closed;

\item if \(\nu\in\Delta\) satisfies
\[
\nu P=\rho_\nu\nu
\qquad\text{with}\qquad
\rho_\nu\neq\rho_\kappa,
\]
then
\[
\nu h=0
\qquad\text{and}\qquad
\nu(U)=0;
\]

\item for every \(i\in U\), there exists an index
\(\beta(i)\in\{1,\ldots,L\}\) such that \(S_{\beta(i)}\subset U\),
the class \(S_{\beta(i)}\) is accessible from \(i\) along a finite
\(P\)-path contained in \(U\), and \(S_{\beta(i)}\) reaches no distinct
communicating class contained in \(U\);

\item the class \(S_{\beta(i)}\) in \textup{(iii)} satisfies
\[
\rho_{\beta(i)}=\rho_\kappa,
\qquad
\min_{x\in S_{\beta(i)}}h(x)>0;
\]

\item every \(i\in U\) can reach \(S_\kappa\).
\end{enumerate}
\end{lemma}

\begin{hypothesis}[Initial support]
\label{hyp:initial-support}
The initial condition satisfies, almost surely,
\[
S=\overline{\supp(\mu_0)}
\qquad\text{and}\qquad
X_0\in\supp(\mu_0).
\]
\end{hypothesis}

\begin{lemma}[Comparison of weights]
\label{lem:controlled-window-weights}
For every \(\Theta>0\), there exist \(c_\Theta>0\) and
\(N_\Theta\geq n_0\) such that, whenever \(N_\Theta\leq m\leq n\) and
\(\sum\limits_{j=m+1}^{n}\gamma_j\leq\Theta\), one has
\[
\frac{w_m}{w_n}\geq c_\Theta.
\]
\end{lemma}

\begin{lemma}[Block-time limsup estimate]
\label{lem:block-time-limsup}
With the notation above,
\[
\limsup_{n\to\infty}
\frac{\mu_{T_n}h}{\gamma_{T_n}}
=
+\infty
\qquad\text{almost surely}.
\]
\end{lemma}

\begin{proof}
Set
\[
R_n:=\frac{\mu_nh}{r_n}.
\]
By the weighted occupation formula,
\[
R_n
=
w_0\mu_0h+\sum_{k=1}^n w_kh(X_k),
\qquad
\frac{\mu_nh}{\gamma_n}
=
\frac{R_n}{w_n}.
\]
In particular, \((R_n)\) is non-decreasing. Since \(U^c\) is closed, Hypothesis~\ref{hyp:initial-support} gives
\(\mu_0(U)>0\) almost surely. Hence the events
\[
E_i:=\{\mu_0(i)>0\},
\qquad i\in U,
\]
cover an almost sure event. Fix \(i\in U\) and work on \(E_i\).


By Lemma~\ref{lem:right-perron-support}\textup{(iii)--(iv)}, there exists an
index \(\beta(i)\) such that \(S_{\beta(i)}\subset U\), the
class \(S_{\beta(i)}\) is accessible from \(i\) along a \(P\)-path contained
in \(U\), and
\[
\rho_{\beta(i)}=\rho_\kappa,
\qquad
h_i^*:=\min_{x\in S_{\beta(i)}}h(x)>0.
\]
Since \(P_{S_{\beta(i)}}\) is irreducible and has positive Perron value, every
state of \(S_{\beta(i)}\) has a successor in \(S_{\beta(i)}\) with positive
transition probability. A path that has entered \(S_{\beta(i)}\) can therefore
be prolonged inside this class for any prescribed finite number of steps.

Fix \(L\geq1\). We may consequently choose
\begin{align}\label{eq:i-path-L-visits}
    i=y_0^{(L)},y_1^{(L)},\ldots,y_{\ell_L}^{(L)}
\end{align}
such that
\begin{align}\label{eq:lambda^L-visits}
\lambda^{(L)}
:=
\prod_{r=0}^{\ell_L-1}
P(y_r^{(L)},y_{r+1}^{(L)})
>0,
\qquad
\sum_{r=1}^{\ell_L}
\mathbf 1_{\{y_r^{(L)}\in S_{\beta(i)}\}}
\geq L.
\end{align}


Now for every \(x\in S\), Hypothesis~\ref{hyp:main} provides an integer \(\ell_x\geq0\) and states
\[
z_0^x=x,z_1^x,\ldots,z_{\ell_x}^x
\]
such that
\[
d_x
:=
\left(
\prod_{r=0}^{\ell_x-1}
P(z_r^x,z_{r+1}^x)
\right)
q(z_{\ell_x}^x)
>0.
\]
By finiteness of \(S\), let
\[
p:=1+\max_{x\in S}\ell_x,
\qquad
d:=\min_{x\in S}d_x>0.
\]

For \(n\geq0\), we next define 
\[
\sigma_n:=n+\ell_{X_n}+1. 
\]
Since \(X_n\) is \(\mathcal F_n\)-measurable, \(\sigma_n\) is an
\((\mathcal F_k)\)-stopping time determined at time \(n\), with
\[
n+1\leq\sigma_n\leq n+p.
\]
Let \(\mathcal I_n\in\mathcal F_{\sigma_n}\) be the event that the process follows the path associated with \(X_n\) and then jumps to \(i\):
\[
\mathcal I_n
=
\left\{
X_{n+r}=z_r^{X_n},\ 1\leq r\leq\ell_{X_n},
\quad
X_{\sigma_n}=i
\right\}.
\]
% Then \(\mathcal I_n\in\mathcal F_{\sigma_n}\). 
Observe that along the prescribed path \(K_\mu\geq P\), while at its endpoint
\[
K_{\mu_{n+\ell_{X_n}}}
\bigl(z_{\ell_{X_n}}^{X_n},i\bigr)
\geq
q(z_{\ell_{X_n}}^{X_n})
\mu_{n+\ell_{X_n}}(i).
\]
Moreover, since \(\ell_{X_n}\leq p-1\) and \((r_n)\) is non-increasing,
\[
\mu_{n+\ell_{X_n}}(i)
\geq
r_{n+\ell_{X_n}}\mu_0(i)
\geq
r_{n+p}\mu_0(i).
\]
Therefore
\[
\mathbb P(\mathcal I_n\mid\mathcal F_n)
\geq
d\,\mu_0(i)r_{n+p}.
\]

We now consider the auxiliary deterministic grid
\[
u_q:=qb_L,
\quad q\geq0, \qquad \text{ where } \qquad b_L:=p+\ell_L
\]
and define \(\mathcal H_q:=\mathcal F_{u_q}\) for \(q\geq0\). The number \(b_L\) is a uniform upper bound on the duration of the prescribed
excursion; it need not be its exact duration.
The choice of \(b_L\) ensures that, irrespective of \(X_{u_q}\),
\[
\sigma_{u_q}+\ell_L
\leq
u_q+p+\ell_L
=
u_{q+1}.
\]
Thus the two prescribed parts of the excursion are completed by time
\(u_{q+1}\).



Let
\[
\mathcal J_q^{(L)}
:=
\left\{
X_{\sigma_{u_q}+r}=y_r^{(L)},
\ 1\leq r\leq\ell_L
\right\},
\]
and set
\[
\mathcal R_q^{(L)}
:=
\mathcal I_{u_q}\cap\mathcal J_q^{(L)}.
\]
On \(\mathcal R_q^{(L)}\), the process first reaches \(i\) through the path
associated with \(X_{u_q}\), and then follows \(y^{(L)}\). Since
\(\sigma_{u_q}+\ell_L\leq u_{q+1}\),
\[
\mathcal R_q^{(L)}\in\mathcal F_{u_{q+1}}.
\]



On \(\mathcal I_{u_q}\), one has \(X_{\sigma_{u_q}}=i\). By successive
conditioning and the inequality \(K_\mu\geq P\),
\[
\mathbb P\left(
\mathcal J_q^{(L)}
\,\middle|\,
\mathcal F_{\sigma_{u_q}}
\right)
\geq
\lambda^{(L)}
\qquad\text{on }\mathcal I_{u_q}.
\]
Since \(\mathcal I_{u_q}\in\mathcal F_{\sigma_{u_q}}\), the tower property
gives
\[
\begin{aligned}
\mathbb P\left(
\mathcal R_q^{(L)}
\,\middle|\,
\mathcal F_{u_q}
\right)
&=
\mathbb E\left[
\mathbf1_{\mathcal I_{u_q}}
\mathbb P\left(
\mathcal J_q^{(L)}
\,\middle|\,
\mathcal F_{\sigma_{u_q}}
\right)
\,\middle|\,
\mathcal F_{u_q}
\right]\\
&\geq
\lambda^{(L)}
\mathbb P\left(
\mathcal I_{u_q}
\,\middle|\,
\mathcal F_{u_q}
\right)\\
&\geq
d\lambda^{(L)}\mu_0(i)r_{u_q+p}.
\end{aligned}
\]
Set
\[
a_L:=d\lambda^{(L)}\mu_0(i).
\]
The random variable \(a_L\) is \(\mathcal F_0\)-measurable and strictly
positive on \(E_i\).
Since \(u_q+p\leq u_{q+1}\),
\[
\mathbb P\left(
\mathcal R_q^{(L)}
\,\middle|\,
\mathcal F_{u_q}
\right)
\geq
a_Lr_{u_{q+1}}
\qquad\text{on }E_i.
\]
The intervals \([u_q,u_{q+1})\) partition \(\mathbb N\), and
\[
\sum_{n=u_q}^{u_{q+1}-1}r_n
\leq
b_Lr_{u_q}.
\]
By the weak reinforcement regime condition \(\sum_n r_n=\infty\),
\[
\sum_qr_{u_q}=\infty,
\qquad
\sum_qr_{u_{q+1}}=\infty.
\]
Thus, on \(E_i\),
\[
\sum_{q\geq0}
\mathbb P\left(
\mathcal R_q^{(L)}
\,\middle|\,
\mathcal F_{u_q}
\right)
=
\infty.
\]
Since
\(\mathcal R_q^{(L)}\in\mathcal H_{q+1}\),
Lévy's conditional Borel--Cantelli lemma gives
\[
\mathbb P\left(
E_i\cap
\{\mathcal R_q^{(L)}\text{ infinitely often}\}
\right)
=
\mathbb P(E_i).
\]


On \(\mathcal R_q^{(L)}\), let
\[
\Lambda_q^{(L)}
:=
\left\{
\sigma_{u_q}+r:
1\leq r\leq\ell_L,\
y_r^{(L)}\in S_{\beta(i)}
\right\},
\]
such that on \(\mathcal R_q^{(L)}\),
\[
|\Lambda_q^{(L)}|\geq L,
\qquad
h(X_t)\geq h_i^*
\quad\text{for }t\in\Lambda_q^{(L)}.
\]

Let
\[
N_q:=\inf\{n:T_n\geq u_{q+1}\}.
\]
Then \(N_q\to\infty\), and, for all sufficiently large \(q\),
\[
T_{N_q-1}<u_{q+1}\leq T_{N_q}.
\]
For \(t\in\{u_q,\ldots,u_{q+1}\}\),
\[
\sum_{j=t+1}^{T_{N_q}}\gamma_j
\leq
\sum_{j=u_q+1}^{u_{q+1}}\gamma_j
+
\sum_{j=u_{q+1}+1}^{T_{N_q}}\gamma_j.
\]
For all sufficiently large \(q\),
\[
\sum_{j=u_q+1}^{u_{q+1}}\gamma_j
\leq
b_L\gamma_{u_q+1}
\leq1,
\]
while
\[
\sum_{j=u_{q+1}+1}^{T_{N_q}}\gamma_j
<
\tau+\gamma_{T_{N_q}}
\leq\tau+1.
\]
Therefore
\[
\sum_{j=t+1}^{T_{N_q}}\gamma_j\leq\tau+2.
\]
For all sufficiently large \(q\), with
\(c_\tau:=c_{\tau+2}\), Lemma~\ref{lem:controlled-window-weights} gives
\[
\frac{w_t}{w_{T_{N_q}}}
\geq
c_\tau,
\qquad
u_q\leq t\leq u_{q+1}.
\]

Notice that the prescribed excursion need not occupy the whole interval \([u_q,u_{q+1}]\). If it is completed earlier, its contributions have already
been added to the non-decreasing quantity \(R_n\). The preceding comparison shows that their weights remain uniformly comparable to
\(w_{T_{N_q}}\).



Consequently, on \(\mathcal R_q^{(L)}\),
\[
\frac{\mu_{T_{N_q}}h}{\gamma_{T_{N_q}}}
=
\frac{R_{T_{N_q}}}{w_{T_{N_q}}}
\geq
\sum_{t\in\Lambda_q^{(L)}}
\frac{w_t}{w_{T_{N_q}}}h(X_t)
\geq
c_\tau Lh_i^*.
\]
Since \(\mathcal R_q^{(L)}\) occurs infinitely often almost surely on \(E_i\)
and \(N_q\to\infty\), it follows that, for this fixed \(L\),
\[
\limsup_{n\to\infty}
\frac{\mu_{T_n}h}{\gamma_{T_n}}
\geq
c_\tau Lh_i^*
\qquad\text{almost surely on }E_i.
\]

Taking the intersection over \(L\in\mathbb N^*\), and recalling that \(c_\tau h_i^*>0\) is independent of \(L\), we obtain
\[
\limsup_{n\to\infty}
\frac{\mu_{T_n}h}{\gamma_{T_n}}
=
+\infty
\qquad\text{almost surely on }E_i.
\]
Finally, since \(U\) is finite and the events \(E_i\), \(i\in U\), cover an almost sure event, the result follows.



\end{proof}

\end{document}
